In this section, we will consider one application of the IQC duality result,
namely robust controller gain synthesis, i.e.
$u(t) = \mcl{K}\mbf{x}_f(t)$. We will represent the corresponding PIE system in
the generalized plant format.

The IQC in \eqref{eqn:unciqc} can be written as
\[
	\ip{p_T {\bmat{v_1\\v_2}}}
	{p_T \mcl{M}_{\mcl{V}_\Delta}
		\bmat{v_1\\v_2}}_{\mbf L_{[a,b]}} \ge 0,
\]
where \(\bmat{v_1\\v_2}:=v:=\Psi\bmat{z_\Delta \\ w_\Delta}\) is the output of the linear system
\(\Psi\) driven by \(z_\Delta\) and \(w_\Delta\) with zero initial conditions.
Let \(\Psi\) admit the realization
\begin{equation}
	\bmat{\mcl{T}_\Psi \dot{\mbf x}_\Psi(t) \\ v_1(t)\\v_2(t)}
	=
	\bmat{\mcl{A}_\Psi & \mcl{B}_\Psi^{z_\Delta} & \mcl{B}_\Psi^{w_\Delta} \\
		\mcl{C}_{\Psi,1} & \mcl{D}_{\Psi,1}^{z_\Delta} & \mcl{D}_{\Psi,1}^{w_\Delta}\\
		\mcl{C}_{\Psi,2} & \mcl{D}_{\Psi,2}^{z_\Delta} & \mcl{D}_{\Psi,2}^{w_\Delta}}
	\bmat{\mbf x_\Psi(t) \\ z_\Delta(t) \\ w_\Delta(t)}.
\end{equation}
For the basis $\Psi = \bmat{\Psi_{11} & \Psi_{12} \\ \Psi_{21} & \Psi_{22}}$ function we require $\Psi_{22}$ invertible,
such that the Potapov--Grinsburg transform of \(\Psi\) is well-defined. Typically, for the dynamic multipliers it results in the form
$\Psi =\bmat{\Psi_{v} & \Psi_{v}\\ \Psi_{v} & \Psi_{v}}$, i.e. if a partition is populated every partition has the same dynamics.
Typically $\Psi_{v}$ is taken tall such that it can represent the spectral landscape of the IQC well. We resolve this issue by 
'squaring up' the dynamic multiplier. If $\Psi_{v} = \bmat{F_0(s) \\ F_1(s) \\ \vdots \\ F_v(s)}$, then this
can be squared up by taking $\Psi_\square = \bmat{F_0(s) & 0\\ F_{1\to v}(s) & I}$, this can be inverted if $F_0(s)$ is invertible,
then $\Psi_\square^{-1} = \bmat{F_0^{-1}(s) & 0\\ -F_{1\to v}(s)F_0(s)^{-1} & I}$. This squaring up also introduces a lifted statespace.
\begin{definition}\label{def:uPIE}[Generalized Plant]
	Let the PI operators be
	$\mathcal{T}, \mathcal{A}: \mathcal{L}(\RL^{\mbf{n}_x}), \quad
		\mathcal{B}_\Delta: \RL^{\mbf{n}_{w_\Delta}} \to \RL^{\mbf{n}_x}, \quad
		\mathcal{B}_w: \R^{n_w} \to \RL^{\mbf{n}_x}, \quad
		\mathcal{C}_\Delta: \RL^{\mbf{n}_x} \to \RL^{\mbf{n}_{z_\Delta}}, \quad
		\mathcal{D}_\Delta: \R^{\mbf{n}_{w_\Delta}} \to \RL^{\mbf{n}_{z_\Delta}}, \quad
		\mathcal{D}_{\Delta w}: \R^{n_{w}} \to \RL^{\mbf{n}_{z_\Delta}}, \quad ,
		\mathcal{C}_z: \RL^{\mbf{n}_{x}} \to \R^{n_z}, \quad
		\mathcal{D}_{z\Delta}: \RL^{\mbf{n}_{w_\Delta}} \to \R^{n_{z}}, \quad ,
		\mathcal{D}_{zw} \in \R^{n_z \times n_w}$,
	with $\mathbf{x}_f \in \RL^{\mbf{n}_x}$, regulated output
	$z\in L_{2,e}^{n_z}[0,\infty)$, disturbance $w\in L_2^{n_w}[0,\infty)$, and uncertainty $\Delta \in \mathbf{\Delta}$, defined according to Definition~\ref{def1}, construct the following uncertain PIE,
	\begin{align}
		\mathcal{T}\dot{\mathbf{x}}_f(t) & = \tilde{\mathcal{A}}(\delta)\mathbf{x}_f(t)+\tilde{\mathcal{B}}_w(\delta)w(t),\notag \\
		z(t)                             & = \tilde{\mathcal{C}}_z(\delta)\mathbf{x}_f(t)+\tilde{D}_{zw}(\delta)w(t).
		\label{eq:uPIE}
	\end{align}
	The uncertain PIE \eqref{eq:uPIE} is then defined to admit a $[G, \Delta]$
	form according to Definition \ref{def1} if there exist PI operators
	$\{\mcl T, \mcl A, \mcl B_\Delta, \mcl B_w, \mcl C_\Delta, \mcl D_\Delta,
		\mcl D_{\Delta w}, \mcl C_z, \mcl D_{z\Delta}, \mcl D_{zw}\}$ that
	parametrize the following PIEs for all
	$(\mathbf{x}_f(t), w_\Delta(t), z_\Delta(t), w(t), z(t))$, and for all
	$\Delta\in\mathbf{\Delta}$,
	\begin{equation}
		{G :\left\{\begin{aligned}
				\mathcal{T}\dot{{\mathbf{x}}}_f(t) & = \mathcal{A}{\mathbf{x}}_f(t) + \mathcal{B}_\Delta w_\Delta(t) +\mathcal{B}_w w(t),\notag                \\
				z_\Delta(t)                        & = \mathcal{C}_\Delta\mathbf{x}_f(t) + \mathcal{D}_\Delta w_\Delta(t) + \mathcal{D}_{\Delta w} w(t),\notag \\
				z(t)                               & = \mathcal{C}_z\mathbf{x}_f(t) + \mathcal{D}_{z\Delta} w_\Delta(t) + \mathcal{D}_{zw} w(t),
			\end{aligned}\right.}
	\end{equation}
	\begin{center}
		with, $w_\Delta(t) = (\Delta z_\Delta)(t).$
	\end{center}
\end{definition}

This definition is adapted from Corollay 8 in \cite{10156335}
\begin{definition}[Primal and Dual Augmented Systems]\label{def:partitioned_augmented_systems}
	\noindent \textbf{Primal System \( G \)}:
	The primal system \( G \) is defined by the following equations:
	The augmented realization uses the partitioned filter
	output \(\bmat{v_1\\v_2}=\Psi\bmat{z_\Delta\\w_\Delta}\). We assume
	there is a bounded selection operator \(\mcl N\) such that
	\(w_\Delta(t)=\mcl N v_2(t)\). The augmented input--output map is then
	written from \(\bmat{v_2^*&w^*}^*\) to \(\bmat{v_1^*&z^*}^*\).
	\begin{equation}
		\begin{aligned}
			\bmat{\mcl{T} & 0 \\ 0 & \mcl{T}_{\Psi}}
			\bmat{{\dot{\mbf{x}}}_f(t) \\ \dot{\mbf{x}}_\Psi(t)}
				&=
			\bmat{
				\mcl{A} & 0 \\
				\mcl{B}_{\Psi}^{z_\Delta}\mcl{C}_{\Delta} & \mcl{A}_\Psi
			}
			\bmat{{\mbf{x}}_f(t) \\ \mbf{x}_\Psi(t)} \\
				&\quad+
			\bmat{
				\mcl{B}_{\Delta}\mcl{N} & \mcl{B}_w \\
				\mcl{B}^{z_\Delta}_\Psi\mcl{D}_\Delta\mcl{N} + \mcl{B}_\Psi^{w_\Delta}\mcl{N} &
				\mcl{B}_{\Psi}^{z_\Delta}\mcl{D}_{\Delta w}
			}
			\bmat{v_2(t) \\ w(t)}, \\
			{\bmat{v_1(t)\\z(t)}}
				&=
			{\bmat{
				\mcl{D}^{z_\Delta}_{\Psi,1}\mcl{C}_\Delta & \mcl{C}_{\Psi,1} \\
				\mcl{C}_z & 0
			}}
			\bmat{{\mbf{x}}_f(t) \\ \mbf{x}_\Psi(t)} \\
				&\quad+
			{\bmat{
				\mcl{D}^{z_\Delta}_{\Psi,1}\mcl{D}_\Delta\mcl{N} + \mcl{D}^{w_\Delta}_{\Psi,1}\mcl{N} &
				\mcl{D}^{z_\Delta}_{\Psi,1}\mcl{D}_{\Delta w} \\
				\mcl{D}_{z\Delta}\mcl{N} & \mcl{D}_{zw}
			}}
			\bmat{v_2(t) \\ w(t)}.
		\end{aligned}
	\end{equation}
\end{definition}
Say we want to control the given PIE system
\begin{equation*}
	\bmat{\partial_t(\mcl{T}\mbf{x}_f(t)) \\ z(t)}
	=
	\bmat{
		\bar{\mcl{A}}(\delta) & \bar{\mcl{B}}_w(\delta) & \bar{\mcl{B}}_u(\delta) \\
		\bar{\mcl{C}}(\delta) & \bar{\mcl{D}}_{zw}(\delta) & \bar{\mcl{D}}_{zu}(\delta)
	}
	\bmat{\mbf{x}_f(t) \\ w(t) \\ u(t)}
\end{equation*}
That admits the generalized plant
$\{\mcl{T}, (\mcl{A}+\mcl{B}_u\mcl{K}), \mcl{B}_{\Delta}, \mcl{B}_{w},
	(\mcl{C}_\Delta + \mcl{D}_{\Delta u}\mcl{K}), \mcl{D}_\Delta,
	\mcl{D}_{\Delta w}, (\mcl{C}_z+\mcl{D}_{zu}\mcl{K}),
	\mcl{D}_{z\Delta}, \mcl{D}_{zw}\}$, i.e.,
\begin{equation*}
	\bmat{\partial_t(\mcl{T}\mbf{x}_f(t)) \\ z_\Delta(t) \\ z(t)}
	=
	\bmat{
		\mcl{A}+\mcl{B}_u\mcl{K} & \mcl{B}_{\Delta} & \mcl{B}_{w} \\
		\mcl{C}_\Delta + \mcl{D}_{\Delta u}\mcl{K} & \mcl{D}_\Delta & \mcl{D}_{\Delta w} \\
		\mcl{C}_z+\mcl{D}_{zu}\mcl{K} & \mcl{D}_{z\Delta} & \mcl{D}_{zw}
	}
	\bmat{\mbf{x}_f(t) \\ w_\Delta(t) \\ w(t)}
\end{equation*}
with,
\begin{center}
	$w_\Delta(t) = (\Delta z_\Delta)(t).$
\end{center}
The augmented input--output plant associated with
Definition~\ref{def:partitioned_augmented_systems} then becomes
\begin{equation*}
	\begin{aligned}
		\hat{\mcl{T}}
			&= \bmat{\mcl{T} & 0 \\ 0 & \mcl{T}_\Psi}, \\
		\hat{\mcl{A}}
			&= \bmat{
				\mcl{A} & 0 \\
				\mcl{B}_\Psi^{z_\Delta}\mcl{C}_\Delta & \mcl{A}_\Psi
			}, \\
		\hat{\mcl{K}}
			&= \bmat{\mcl{K} & 0 \\ 0 & 0}, \\
		\tilde{\mcl{A}}
			&= \bmat{
				\mcl{B}_u & 0 \\
				\mcl{B}^{z_\Delta}_\Psi\mcl{D}_{\Delta u} & 0
			}, \\
		\hat{\mcl{B}}
			&= \bmat{
				\mcl{B}_{\Delta}\mcl N & \mcl{B}_w \\
				\left(\mcl{B}^{z_\Delta}_\Psi\mcl{D}_\Delta
				+\mcl{B}_\Psi^{w_\Delta}\right)\mcl N &
				\mcl{B}_{\Psi}^{z_\Delta}\mcl{D}_{\Delta w}
			}, \\
		\hat{\mcl{C}}
			&= \bmat{
				\mcl{D}^{z_\Delta}_{\Psi,1}\mcl{C}_\Delta & \mcl{C}_{\Psi,1} \\
				\mcl{C}_z & 0
			}, \\
		\tilde{\mcl{C}}
			&= \bmat{
				\mcl{D}_{\Psi,1}^{z_\Delta}\mcl{D}_{\Delta u} & 0 \\
				\mcl{D}_{zu} & 0
			}, \\
		\hat{\mcl{D}}
			&= \bmat{
				\left(\mcl{D}^{z_\Delta}_{\Psi,1}\mcl{D}_\Delta
				+\mcl{D}^{w_\Delta}_{\Psi,1}\right)\mcl N &
				\mcl{D}^{z_\Delta}_{\Psi,1}\mcl{D}_{\Delta w} \\
				\mcl{D}_{z\Delta}\mcl N & \mcl{D}_{zw}
			}.
	\end{aligned}
\end{equation*}
such that
\begin{equation*}
	\begin{aligned}
		\partial_t(\hat{\mcl{T}}{\hat{\mbf{x}}}(t))
			&= (\hat{\mcl{A}}+\tilde{\mcl{A}}\hat{\mcl{K}})
			\hat{\mbf{x}}(t)+\hat{\mcl{B}}\hat{w}(t), \\
		\hat{z}(t)
			&= (\hat{\mcl{C}}+\tilde{\mcl{C}}\hat{\mcl{K}})
			\hat{\mbf{x}}(t)+\hat{\mcl{D}}\hat{w}(t)
	\end{aligned}
\end{equation*}
where $\hat{\mbf{x}} = \bmat{\mbf{x}_f(t)^* & \mbf{x}_\Psi(t)^*}^*$,
$\hat{z} = \bmat{v_1(t)^* & z(t)^*}^*$,
$\bmat{v_1^*&v_2^*}^* = \Psi\bmat{z_\Delta^*, w_\Delta^*}^*$, and
$\hat{w}(t) = \bmat{v_2(t)^* & w(t)^*}^*$. Then the dual system is
\begin{equation*}
	\begin{aligned}
		\hat{\mcl{T}}^*\dot{\hat{\bar{\mbf{x}}}}(t)
			&= (\hat{\mcl{A}}^*+\hat{\mcl{K}}^*\tilde{\mcl{A}}^*)
			\hat{\bar{\mbf{x}}}(t) \\
			&\quad+ (\hat{\mcl{C}}^*+\hat{\mcl{K}}^*\tilde{\mcl{C}}^*)
			\hat{\bar{w}}(t), \\
		\hat{\bar{z}}(t)
			&= \hat{\mcl{B}}^*\hat{\bar{\mbf{x}}}(t)
			+ \hat{\mcl{D}}^*\hat{\bar{w}}(t)
	\end{aligned}
\end{equation*}
Moreover we also define
\begin{equation}\label{eqn:PNPIMult}
	\tilde{\mcl{V}}
	=
	\bmat{
		\tilde{\mcl{V}}_{11} & \tilde{\mcl{V}}_{12} \\
		\tilde{\mcl{V}}_{21} & \tilde{\mcl{V}}_{22}
	}
	=
	\bmat{
		\mcl{V}_{11} & 0 & \mcl{V}_{12} & 0 \\
		0 & \hat{\mcl{V}}_{11} & 0 & \hat{\mcl{V}}_{12} \\
		\mcl{V}_{21} & 0 & \mcl{V}_{22} & 0 \\
		0 & \hat{\mcl{V}}_{21} & 0 & \hat{\mcl{V}}_{22}
	}
\end{equation}
Here the first and third block rows/columns of \(\tilde{\mcl V}\) correspond
to the partition \(v=\bmat{v_1^* & v_2^*}^*\). Since
\(w_\Delta=\mcl N v_2\), the static multiplier is applied to the genuine
output--input pair of the augmented map, and \eqref{eqn:PNPIMult} acts on
\(\bmat{v_1^*&z^*&v_2^*&w^*}^*\).
If $\mcl{V}$ and $\hat{\mcl{V}}$ are {strict PN-PI} operators, then $\tilde{\mcl{V}}$ is also a {strict PN-PI} operator. If we would derive \textbf{robust performance} LPI without using the duality
theorem, we would end up with a non-convex LPI, which we cannot efficiently
solve using the PIETOOLS library. Therefore we apply the duality theorem to
derive the following LPI.

\begin{algorithm}[t]
	\caption{Programming the Dual Synthesis LPI}
	\label{alg:program_dual_lpi}
	\begin{algorithmic}[1]
		\Require Generalized plant operators, basis degree \(\nu\), pole \(\rho\),
		and multiplier templates for \(\mcl V_\Delta\) and \(\hat{\mcl V}\).
		\State Construct the basis filter and partition
		\(v_1=(\psi_\nu\otimes I_{z_\Delta})z_\Delta\),
		\(v_2=(\psi_\nu\otimes I_{w_\Delta})w_\Delta\).
		\State Set \(\mcl N=\bmat{I_{w_\Delta}&0&\cdots&0}\), so that
		\(w_\Delta=\mcl N v_2\).
		\State Substitute \(w_\Delta=\mcl N v_2\) and form the augmented map
		\(\hat G:\bmat{v_2^*&w^*}^*\mapsto\bmat{v_1^*&z^*}^*\).
		\State Assemble the fixed operators
		\(\hat{\mcl T},\hat{\mcl A},\hat{\mcl B},\hat{\mcl C},\hat{\mcl D},
		\tilde{\mcl A},\tilde{\mcl C}\).
		\State Declare decision variables \(\mcl P=\mcl P^*\succ0\),
		\(\mcl Z=\bmat{\mcl Z_{11}&\mcl Z_{12}\\0&0}\), and strict PN-PI
		multiplier variables.
		\State Build \(\tilde{\mcl V}\) from \eqref{eqn:PNPIMult} and compute its
		dual multiplier \(\tilde{\mcl V}'\).
		\State Assemble the LPI in Theorem~\ref{cor:robSYN}, using
		\(\mcl Z=\hat{\mcl K}\mcl P\) instead of the controller
		\(\hat{\mcl K}\).
		\State Solve the LPI feasibility problem.
		\State Recover the controller as
		\(\mcl K=\mcl Z_{11}\mcl P_{11}^{-1}\).
	\end{algorithmic}
\end{algorithm}

\begin{theorem}[Robust Controller Gain Synthesis]\label{cor:robSYN}

	Let $G$ be parametrized by PI operators
	$\{\mcl{T}, (\mcl{A}+\mcl{B}_u\mcl{K}), \mcl{B}_{\Delta}, \mcl{B}_{w},
		(\mcl{C}_\Delta + \mcl{D}_{\Delta u}\mcl{K}), \mcl{D}_\Delta,
		\mcl{D}_{\Delta w}, (\mcl{C}_z+\mcl{D}_{zu}\mcl{K}),
		\mcl{D}_{z\Delta}, \mcl{D}_{zw}\}$. Let $\Delta\in\mbf{\Delta}$, and
	there exists a {strict PN-PI} multiplier $\mcl{M}_{\tilde{\mcl{V}}}$ parameterized
	according to \eqref{eqn:PNPIMult}, such that

	\begin{itemize}
		\item The interconnection $[G,\Delta]$ is well--posed for all $\Delta\in\mbf{\Delta}$
		\item The IQC in Definition \ref{def:multipliers} is satisfied for all $\Delta\in\mbf{\Delta}$.
		\item The dynamic multiplier admits the partition
		\(v=\bmat{v_1^*&v_2^*}^*\) and a bounded selection operator \(\mcl N\)
		such that \(w_\Delta(t)=\mcl N v_2(t)\).
		\item The hatted operators
		\(\{\hat{\mcl T},\hat{\mcl A},\hat{\mcl B},\hat{\mcl C},\hat{\mcl D}\}\)
		denote the causal bounded realization
		\(\hat G:\bmat{v_2^*&w^*}^*\mapsto\bmat{v_1^*&z^*}^*\) constructed
		above, and this realization and its dual scaled realization satisfy Theorem~\ref{thm:dualscaledIO}.
	\end{itemize}
	Then the interconnection $[G,\Delta]$ is \textbf{robustly stable}, with
	$\hat{\mathbf{x}}(0) = 0$, and \textbf{robust performance} is achieved if
	there exist a PI operator
	$\mcl{Z}: {\mcl{L}\left( \mathbb{R}^{\mbf{n}_{\hat{\mathbf{x}}}}
		\to \mathbb{R}^{2n_u} \right)}$, such that
	$\mcl{Z} = \bmat{\mcl{Z}_{11} & \mcl{Z}_{12} \\ 0 & 0}
		= \bmat{\mcl{K}\mcl{P}_{11} & \mcl{K}\mcl{P}_{12} \\ 0 & 0}$,
	a strict PN-PI multiplier $\mcl{M}_{\tilde{\mcl{V}}}$, with
	$\hat{\mcl V}_{11} \succeq 0$, and a coercive self-adjoint
	$\mathcal{P}\in \Pi_4$ with ${\mathcal{P}=\mathcal{P}^*\succ 0}$, such that,

	\begin{itemize}
		\item $\mathcal{K} = \mathcal{Z}_{11}\left(\mathcal{P}_{11}\right)^{-1}$ ,

		\item and
		      \vspace{-2ex}
		      \begin{equation*}
			      \bmat{
				      \hat{\mcl{A}}\mcl{P}\hat{\mcl{T}}^* + \tilde{\mcl{A}}\mcl{Z}\hat{\mcl{T}}^* + (*)
				      & \hat{\mcl T}\mcl P\hat{\mcl C}^*
				      + \hat{\mcl T}\mcl Z^*\tilde{\mcl C}^*
				      + \hat{\mcl B}\tilde{\mcl V}'_{12}
				      & \hat{\mcl{B}}(\tilde{\mcl{V}}'_{11})^* \\
				      \hat{\mcl{C}}\mcl{P}\hat{\mcl{T}}^* + \tilde{\mcl{C}}\mcl{Z}\hat{\mcl{T}}^* + \tilde{\mcl{V}}'_{21}\hat{\mcl{B}}^*
				      & \tilde{\mcl{V}}'_{21}\hat{\mcl{D}}^* + \hat{\mcl{D}}\tilde{\mcl{V}}'_{12} + \tilde{\mcl{V}}'_{22} + \epsilon I
				      & \hat{\mcl{D}}(\tilde{\mcl{V}}'_{11})^* \\
				      \tilde{\mcl{V}}'_{11}\hat{\mcl{B}}^*
				      & \tilde{\mcl{V}}'_{11}\hat{\mcl{D}}^*
				      & -\tilde{\mcl{V}}'_{11}
			      }
			      \preceq 0
		      \end{equation*}
	\end{itemize}
	Here, $\tilde{\mathcal{V}}'$ is the dual multiplier of $\tilde{\mcl{V}}$
\end{theorem}
\begin{proof}

	Suppose the LPI holds. Define
	$V(\hat{\bar{\mbf{x}}}) =
		{\langle\hat{\mcl{T}}^*\hat{\bar{\mbf{x}}},
		\mcl{P}\hat{\mcl{T}}^*\hat{\bar{\mbf{x}}}\rangle} \geq 0$.
	{The LPI is derived in Appendix~\ref{app:robSYN_LPI_derivation}.}
	We have that,
	\begin{align*}
		\dot{V}(\hat{\bar{\mbf{x}}}(t))
			&{+ \epsilon\|\hat{\bar w}(t)\|^2
				+ \left\langle
				\bmat{\hat{\bar z}(t)\\ \hat{\bar w}(t)},
				\tilde{\mcl V}'\bmat{\hat{\bar z}(t)\\ \hat{\bar w}(t)}
				\right\rangle_{\RL}}
			\leq 0
	\end{align*}
	Now, since $V(\hat{\bar{\mbf{x}}}(0)) = V(0) = 0$ and
	$V(\hat{\bar{\mbf{x}}}(T))\geq0$, after integrating in time, we obtain,
	\begin{equation*}
		\begin{aligned}
			&{\int_0^T \left\langle
				\bmat{\hat{\bar z}(t)\\ \hat{\bar w}(t)},
				\tilde{\mcl V}'\bmat{\hat{\bar z}(t)\\ \hat{\bar w}(t)}
			\right\rangle_{\RL}dt} \leq {-\epsilon\int_0^T\|\hat{\bar w}(t)\|^2dt}
		\end{aligned}
	\end{equation*}
	Then by Theorem \ref{thm:dualscaledIO}, Lemma \ref{lem:PNPIfact} and
	Lemma \ref{lem:primdualPNPI}, {there exists $\epsilon>0$ such that}
	\begin{equation*}
		\begin{aligned}
			&\int_0^T \left\langle
				\bmat{v_1(t)\\z(t)\\v_2(t)\\w(t)},
				\tilde{\mcl V}\bmat{v_1(t)\\z(t)\\v_2(t)\\w(t)}
			\right\rangle_{\mathbb{R}L_2}dt \\
			&\leq -{\epsilon}
			(\|{{w}_\Delta}(t)\|^2_{\RL}+\| {w}(t)\|^2_{\R})
		\end{aligned}
	\end{equation*}

	Thus we can conclude
	\begin{equation*}
		\begin{gathered}
			\ip{
				p_T\bmat{v_1\\z\\v_2\\w}
			}{
				p_T\mcl M_{\tilde{\mcl V}}\bmat{v_1\\z\\v_2\\w}
			}_{\mbf L_{[a, b]}} \\
			\leq -{\epsilon}
			(\|p_T{w_\Delta}\|^2_{\mbf L_{[a, b]}}+\|p_T w\|^2_{L_2}),
		\end{gathered}
	\end{equation*}
	Therefore, by Theorem 1 in
	\cite{lenssen2025muanalysissynthesisframeworkpartial}, robust performance
	and stability are guaranteed.
\end{proof}
